\chapter{Results and discussion}\label{ch:results}

\section{Model accuracy}
The twin is only as good as its cell model. The OCV curve was rebuilt from the GITT tests of the M50T data set and validated on cells that were not used to build it: the voltage error fell from 64.7\,mV with a literature curve for the sister cell \citep{chen2020} to 17.0\,mV (\cref{fig:ocv}). The temperature dependence of the resistances and the resistance growth with ageing ($R/R_\mathrm{new} = 1 + 1.31\,(1-\mathrm{SOH})$, 40 cells) were also fitted to the measurements.

\fig{ocv_dogrulama}{0.85}{OCV validation on cells not used for fitting.}{fig:ocv}

\section{SOC estimation accuracy}
\Cref{tab:soc} compares SOC estimators on the same reference drive. The EKF with temperature-dependent resistance and offset state has the lowest error.

\begin{table}[htbp]\centering\small
\caption{SOC estimation error on the reference drive (after 60\,s).}\label{tab:soc}
\begin{tabular}{@{}lrr@{}}\toprule
Method & RMSE (\%) & Max.\ error (\%) \\\midrule
OCV lookup & 17.0 & 36.2 \\
Coulomb counting & 0.86 & 1.14 \\
EKF, fixed resistance & 0.72 & 1.09 \\
\textbf{EKF, $R(T)$ + offset (proposed)} & \textbf{0.28} & \textbf{0.52} \\\bottomrule
\end{tabular}
\end{table}

\fig{estimator_comparison_normal}{0.85}{SOC error of the estimators.}{fig:soc}

The start-up offset check detects a current-sensor offset of 4\,A or more within 0.9\,s; without it the filter needed more than 1000\,s or did not detect it at all.

\section{Fault detection and safety}
The twin and the BMS were compiled together and run on the STM32 in 1000 randomised runs of 600\,s each, with random temperature, SOC, driving, cell spread and one fault per run. Every run passed its verdict, and a PC build of the same code gave the same verdict in all 1000 runs (\cref{tab:mc}).

\begin{table}[htbp]\centering\small
\caption{On-target Monte Carlo test, 1000 runs.}\label{tab:mc}
\begin{tabular}{@{}lrl@{}}\toprule
Fault & Runs & Reaction \\\midrule
None & 178 & no false alarm \\
External heat & 151 & detected after 8.0\,s (median) \\
Temperature sensor open & 151 & detected in the same step \\
Current-sensor offset & 202 & detected after 0.9\,s \\
Cooling failure & 199 & detected in 48\,\% (median 138\,s), always safe \\
Aged pack & 119 & resistance growth recognised \\\bottomrule
\end{tabular}
\end{table}

\fig{mc_overview}{0.9}{Monte Carlo results per fault type.}{fig:mc}

On the bench, the hardware tests found three design flaws that simulation had missed: a cooling alarm that switched on and off repeatedly, a false cooling alarm after restarting the BMS on a warm pack, and a core-temperature estimate that lagged far behind the truth under external heat. All three were fixed; with the fix, the safe stop now comes at a true core temperature of 60\,°C instead of 66\,°C.

\section{Cloud synchronisation}
Over 30 minutes the gateway delivered all 18\,046 telemetry rows to the server without loss. The median upload round-trip time was 102\,ms (95th percentile 134\,ms); the end-to-end delay of about 1\,s is mostly the batching interval (\cref{fig:link}).

\fig{link_quality_20261121}{0.85}{Cloud link over 30 minutes.}{fig:link}

\section{Dashboard}
\Cref{fig:dash} shows the dashboard during a bench test. It shows the state and the reason for any current limit, so a user can see why the battery is being protected.

\fig{web_dashboard_live}{0.85}{Web dashboard.}{fig:dash}

\section{Comparison with the base paper}
\Cref{tab:cmp} compares this work with \citet{nejakar2025}. The two SOC errors are not directly comparable: theirs is a mean absolute error on a real 12\,V battery, ours a root-mean-square error against the twin's true SOC.

\begin{table}[htbp]\centering\small
\caption{Comparison with the base paper.}\label{tab:cmp}
\begin{tabularx}{\linewidth}{@{}l>{\raggedright\arraybackslash}X>{\raggedright\arraybackslash}X@{}}\toprule
 & \citet{nejakar2025} & This work \\\midrule
Battery & real 12\,V battery & digital twin of a 19S12P pack (M50T data) \\
Controller & ESP32 (monitoring) & STM32 BMS (monitoring + control) \\
Link & Wi-Fi to Firebase & CAN to BMS, HTTPS to own server \\
SOC method & Kalman + RNN, MAE 2.6\,\% & EKF with offset and resistance states, RMSE 0.28\,\% \\
Alerts & thresholds (SOC $<$ 20\,\%, T $>$ 50\,°C) & four graded states, stress index, anomaly detection \\
Cloud test & 48\,h, no loss, latency $<$ 1\,s & 30\,min, no loss, RTT 102\,ms \\
Validation & real use & 1000 on-target runs + HIL bench \\\bottomrule
\end{tabularx}
\end{table}

\section{Limitations}
The twin sends its measurements over CAN, so the analog measurement hardware of a real BMS is not tested. The BMS does not estimate capacity, so on an aged pack the SOC error rises to about 5\,\%. Results come from a simulated pack; a test on real cells remains future work.
